\section{LEVEL 3 · Steekproefontwerp}

\begin{missie}
Missie: stel een steekproef samen die je kan vertrouwen. (LPD 11)
\end{missie}

\begin{doel}
\textbf{DOEL} \quad Na dit level kan je:

☐ uitleggen wat een \textbf{representatieve} steekproef is en waarom dat belangrijk is;

☐ de \textbf{drie criteria} voor een representatieve steekproef toepassen;

☐ uitleggen wat een \textbf{(enkelvoudig) aselecte steekproef} is;

☐ een \textbf{opportunistische} steekproef en een steekproef op basis van \textbf{vrijwillige respons} herkennen en de gevolgen verklaren.
\end{doel}

\subsection*{Theorie 3.1 · Representatief of niet?}

\begin{theorie}
De manier waarop je een steekproef samenstelt, noemen we het \textbf{steekproefontwerp}. Uit de steekproef moeten we \emph{betrouwbare besluiten} kunnen trekken voor de hele populatie. Daarom moet de steekproef een \textbf{goede weerspiegeling} van de populatie zijn, maar dan met een kleinere omvang: een \textbf{representatieve} steekproef.

Is een steekproef \textbf{niet representatief}, dan wijken de waarden in de steekproef \emph{systematisch} af van de waarden in de populatie. Dan kan je er geen juiste conclusies uit trekken.
\end{theorie}

\begin{theorie}
\textbf{Drie criteria voor een representatieve steekproef}

\begin{enumerate}
\item De \textbf{steekproefomvang} moet voldoende groot zijn.
\item De \textbf{samenstelling} van de steekproef moet overeenkomen met de opbouw van de populatie (bijvoorbeeld leeftijd, geslacht, regio, type organisatie).
\item De \textbf{non-respons} mag niet te groot zijn. Non-respons: geselecteerde personen die niet deelnemen of afhaken.
\end{enumerate}
\end{theorie}

\begin{voorbeeld}
\textbf{Voorbeeld.} Een onderzoeksbureau bevraagt 250 woonzorgcentra over de werkdruk van zorgverleners. De centra worden zo gekozen dat kenmerken die in Vlaanderen vaak voorkomen (regio, grootte, soort centrum) ook in de steekproef vaak voorkomen. Zou de steekproef enkel uit grote commerciële centra in steden bestaan, dan is ze \emph{niet} representatief (criterium 2).

\textbf{Non-respons:} een jeugdbeweging stuurt 120 ouders een vragenlijst; 30 ouders antwoorden. Het responspercentage is $\dfrac{30}{120}\cdot 100\,\%=25\,\%$. Dat is te laag om van een representatieve steekproef te spreken.
\end{voorbeeld}

\subsection*{Theorie 3.2 · Aselecte steekproeven}

\begin{theorie}
\begin{itemize}
\item Een steekproef is \textbf{aselect} als ze op een willekeurige manier, door \textbf{toeval} (randomisatie), wordt samengesteld volgens een vooraf bepaalde werkwijze. Het onderzoeksteam beslist zelf \emph{niet} wie erin komt (niet op basis van ervaring of gezond verstand).
\item Bij een \textbf{enkelvoudig aselecte steekproef (EAS)} heeft \emph{elk} element van de populatie dezelfde kans om in de steekproef te komen, door een volledig willekeurige trekking uit de volledige populatie.
\item Een aselecte steekproef die groot genoeg is en uit de beoogde populatie komt, is zeer waarschijnlijk representatief.
\end{itemize}
\end{theorie}

\begin{voorbeeld}
\textbf{Voorbeeld.} Voor een steekproef van 50 leerlingen uit 800 geef je eerst elke leerling een nummer van 1 tot 800. Daarna laat je de computer lukraak 50 verschillende nummers trekken. Zo is elke leerling even waarschijnlijk: dat is een EAS.
\end{voorbeeld}

\subsection*{Theorie 3.3 · Ontwerpen die tot een vertekende steekproef leiden}

\begin{theorie}
\begin{itemize}
\item \textbf{Opportunistische steekproef}: de onderzoeker kiest elementen die \emph{gemakkelijk, snel of goedkoop} te bereiken zijn. Moeilijker bereikbare elementen hebben dan een kleinere kans. Voorbeeld: enkel mensen bevragen aan de uitgang van het lokale winkelcentrum.
\item \textbf{Steekproef op basis van vrijwillige respons}: personen kiezen zelf of ze deelnemen. Vooral mensen met een \emph{uitgesproken} (zeer positieve of zeer negatieve) mening reageren. Voorbeeld: een online stemming over een geplande fusie.
\end{itemize}
\end{theorie}

\begin{letop}
\textbf{LET OP} \quad Ook bij een representatieve steekproef en een goed onderzoek zal de steekproefwaarde (bijvoorbeeld het gemiddelde) bijna altijd \emph{lichtjes} afwijken van de populatiewaarde. Ze is dan wel \emph{niet systematisch} te groot of te klein. Bij een voldoende grote steekproef blijft de afwijking beperkt.
\end{letop}

\subsection*{Aan de slag}

\subsubsection*{Oefening 3.1 \quad ★☆☆ \quad Invullen}

\textbf{Woordenlijst:} representatief – aselect – non-respons – omvang – samenstelling – vertekend.

Een steekproef die een goede weerspiegeling is van de populatie is \gat{representatief}. Daarvoor moet de \gat{omvang} voldoende groot zijn, moet de \gat{samenstelling} overeenkomen met de populatie en mag de \gat{non-respons} niet te groot zijn. Een \gat{aselect} gekozen steekproef komt door toeval tot stand. Is de steekproef niet representatief, dan is ze \gat{vertekend}.

\subsubsection*{Oefening 3.2 \quad ★★☆ \quad Koppel: begrip en omschrijving}

\breedtes{34/66}
\begin{tabular}{ll}
\hline
Begrip & Omschrijving \\
\hline
1. steekproefontwerp & A. Elk element heeft dezelfde kans om gekozen te worden. \\
2. enkelvoudig aselecte steekproef & B. Mensen kiezen zelf of ze deelnemen; vooral uitgesproken meningen. \\
3. opportunistische steekproef & C. De manier waarop de steekproef wordt samengesteld. \\
4. vrijwillige respons & D. De onderzoeker kiest wie gemakkelijk, snel of goedkoop bereikbaar is. \\
5. non-respons & E. Geselecteerde personen die niet deelnemen of afhaken. \\
\hline
\end{tabular}

Mijn antwoord: 1 – \_\_\quad 2 – \_\_\quad 3 – \_\_\quad 4 – \_\_\quad 5 – \_\_

\begin{oplossing}
\textbf{Oplossing 3.2} \quad 1 – C \quad 2 – A \quad 3 – D \quad 4 – B \quad 5 – E
\end{oplossing}

\subsubsection*{Oefening 3.3 \quad ★★☆ \quad Welk criterium is geschonden?}

Zet een kruisje bij elk criterium dat niet vervuld is: (1) omvang voldoende groot, (2) samenstelling komt overeen met de populatie, (3) non-respons niet te groot.

\breedtes{64/12/12/12}
\begin{tabular}{llll}
\hline
Situatie & 1 & 2 & 3 \\
\hline
a) Een gemeente met 90\,000 inwoners baseert een campagne op een bevraging van 25 inwoners. & & & \\
b) Een onderzoek naar de tevredenheid van ouders gebeurt enkel bij ouders van kinderen in één wijk met veel hoogopgeleide gezinnen. & & & \\
c) Van de 200 uitgenodigde zorgverleners vullen er 14 de vragenlijst in. & & & \\
d) Een bevraging van jongeren gebeurt enkel via een jeugdhuis in een grote stad, bij 40 bezoekers. & & & \\
\hline
\end{tabular}

\begin{oplossing}
\textbf{Oplossing 3.3} \quad a) criterium 1 (25 is veel te weinig). b) criterium 2 (de samenstelling klopt niet). c) criterium 3 (responspercentage $\frac{14}{200}=7\,\%$). d) criteria 1 en 2 (40 is weinig en enkel stedelijke bezoekers).
\end{oplossing}

\subsubsection*{Oefening 3.4 \quad ★★☆ \quad Responspercentage berekenen}

\textbf{a)} Een jeugdbeweging mailt 120 ouders; 30 reageren. Bereken het responspercentage en de non-respons (in \%). \lijn

\textbf{b)} Een woonzorgcentrum nodigt 250 medewerkers uit; 175 doen mee. Bereken het responspercentage. \lijn

\textbf{c)} Welke van de twee bevragingen is het meest betrouwbaar wat non-respons betreft? Verklaar. \lijn

\begin{oplossing}
\textbf{Oplossing 3.4} \quad a) $\frac{30}{120}\cdot 100\,\%=25\,\%$ respons; non-respons 75 \%. b) $\frac{175}{250}\cdot 100\,\%=70\,\%$. c) De bevraging in het woonzorgcentrum: 70 \% respons tegenover 25 \%; er ontbreken veel minder deelnemers die mogelijk anders denken.
\end{oplossing}

\subsubsection*{Oefening 3.5 \quad ★★☆ \quad Is de steekproef aselect?}

\textbf{a)} Een directeur kiest 30 leerlingen ‘die goed kunnen uitleggen’ voor een bevraging over de speelplaats. Is dit aselect? Verklaar. \lijn

\textbf{b)} Beschrijf hoe je een EAS van 40 leerlingen trekt uit een school met 600 leerlingen. \lijn

\textbf{c)} Waarom is een groot genoeg aselecte steekproef waarschijnlijk representatief? \lijn

\begin{oplossing}
\textbf{Oplossing 3.5} \quad a) Nee: de directeur beslist zelf wie meedoet, niet het toeval. b) Elke leerling een nummer van 1 tot 600 geven en met de computer lukraak 40 verschillende nummers trekken. c) Elk element heeft dezelfde kans, dus geen groep is systematisch bevoordeeld; bij voldoende omvang weerspiegelt de steekproef de populatie.
\end{oplossing}

\subsubsection*{Oefening 3.6 \quad ★★☆ \quad Opportunistisch of vrijwillig?}

Welke vorm van een vertekende steekproef is het? Kies ‘opportunistisch’ of ‘vrijwillige respons’ en verklaar kort.

\textbf{a)} Voor een schoolopdracht over pesten vragen leerlingen enkel mensen aan het station of ze meedoen. \lijn

\textbf{b)} Een tijdschrift plaatst een online poll ‘Is de speeltijd te kort?’. \lijn

\textbf{c)} Een sociaal werker bevraagt enkel de bewoners die zij die dag in de leefruimte aantreft. \lijn

\textbf{d)} Een krant vraagt lezers via brief hoe ze vinden dat kinderen opgevoed moeten worden. \lijn

\begin{oplossing}
\textbf{Oplossing 3.6} \quad a) opportunistisch (gemakkelijk bereikbaar). b) vrijwillige respons (wie wil, stemt). c) opportunistisch (enkel wie gemakkelijk bereikbaar is). d) vrijwillige respons (vooral wie een uitgesproken mening heeft, stuurt een brief).
\end{oplossing}

\subsubsection*{Oefening 3.7 \quad ★★★ \quad Context: een steekproef ontwerpen}

Een gemeente met 4\,000 inwoners onder de 18 jaar wil weten welke vrijetijdsactiviteiten jongeren missen.

\textbf{a)} Noteer de populatie, de variabele en het soort variabele. \lijn

\textbf{b)} Beschrijf een steekproefontwerp dat voldoet aan de drie criteria. Geef drie kenmerken waarmee je rekening houdt. \lijn

\textbf{c)} Hoe beperk je de non-respons? \lijn

\begin{oplossing}
\textbf{Oplossing 3.7} \quad (Meerdere antwoorden mogelijk.) a) Populatie: alle 4\,000 inwoners jonger dan 18; variabele: gewenste vrijetijdsactiviteit; kwalitatief, categorisch niet-geordend. b) Aselecte trekking uit het bevolkingsregister van een voldoende grote groep (bv. 400) en rekening houden met leeftijd, geslacht en wijk/deelgemeente. c) Jongeren persoonlijk of meermaals aanschrijven, een korte vragenlijst gebruiken, een herinnering sturen en de bevraging op school laten invullen.
\end{oplossing}

\subsection*{Foutjager}

\textbf{1.} Siska: ‘Ik koos zelf 50 neutrale mensen uit; dat is dus een aselecte steekproef.’ \lijn

\textbf{2.} Karim: ‘Onze online poll kreeg 10\,000 reacties, dus is ze zeker representatief.’ \lijn

\textbf{3.} Emma: ‘Als mijn steekproef representatief is, is het gemiddelde precies gelijk aan dat van de populatie.’ \lijn

\begin{oplossing}
\textbf{Oplossing Foutjager} \quad 1. Zelf kiezen is niet aselect: het toeval moet beslissen. 2. Vrijwillige respons geeft vertekening, ook bij veel reacties. 3. Zelfs een representatieve steekproef wijkt lichtjes af (steekproefvariabiliteit), maar niet systematisch.
\end{oplossing}

\subsection*{CHECK · Level 3}

\textbf{1.} (2 p) Wat is een representatieve steekproef? \lijn

\textbf{2.} (3 p) Noem de drie criteria. \lijn

\textbf{3.} (2 p) Wat is een enkelvoudig aselecte steekproef? \lijn

\textbf{4.} (2 p) Een kinderopvang bevraagt enkel ouders die op de infoavond aanwezig zijn. Welke vorm van vertekening is dit en waarom is dat een probleem? \lijn

\begin{doel}
\textbf{Score:} \_\_\_ / 9 \quad ☐ behaald (vanaf 6 / 9) \quad ☐ nog even oefenen
\end{doel}

\begin{oplossing}
\textbf{Oplossing CHECK 3} \quad 1. Een steekproef die een goede weerspiegeling is van de populatie, maar kleiner. 2. Voldoende grote omvang; samenstelling komt overeen met de populatie; niet te veel non-respons. 3. Elk element van de populatie heeft dezelfde kans, door volledig willekeurige trekking uit de hele populatie. 4. Opportunistisch (en ook vrijwillig): enkel betrokken ouders zijn aanwezig; de anderen ontbreken, dus de resultaten zijn niet te veralgemenen.
\end{oplossing}
